\documentclass[a4paper,twoside]{article}

\usepackage{epsfig}
\usepackage{subcaption}
\usepackage{calc}
\usepackage[hyperfootnotes=false]{hyperref}
\usepackage{xcolor}
\usepackage{amssymb}
\usepackage{amstext}
\usepackage{amsmath}
\usepackage{amsthm}
\usepackage{multicol}
\usepackage{pslatex}
\usepackage{apalike}
\usepackage{algorithm2e}
\usepackage[bottom]{footmisc}
\usepackage{stfloats}
\usepackage{SCITEPRESS}
\usepackage{tabularx}
\usepackage{float}
\usepackage{booktabs}
\usepackage{graphicx}
\usepackage{comment}
\setlength{\emergencystretch}{3em}

% NOTE FOR SUBMISSION: Remove the \author block below for double-blind review.
% NOTE FOR OVERLEAF: Upload step7_lightgbm_fig3_combined.png to the project
%   or compile locally where ../03_outputs/step7/ is accessible.

\begin{document}

\title{Ensemble-Imbalance Machine Learning for ALS Prognosis:\\
A Replication Study Extended with Bayesian Optimisation,\\
LightGBM, and SHAP Interpretability}

\author{\authorname{Ivo Camacho\sup{3}\orcidAuthor{0009-0000-9454-2385},  Pedro Martins\sup{1,2}\orcidAuthor{0000-0000-0000-0000} and Maryam Abbasi\sup{3,4}\orcidAuthor{0000-0002-9011-0734}}
\affiliation{\sup{1}Polytechnic University of Coimbra, Rua da Misericórdia, Lagar dos Cortiços, S. Martinho do Bispo, 3045-093 Coimbra, Portugal}
\affiliation{\sup{3}Polytechnic University of Santarém, Santarém, Portugal.}
\affiliation{\sup{4}Research Hub Digital Literacy \& Social Inclusion (CIAC-PLDIS-IPSantarém), Portugal.}
\email{iscac14461@alumni.iscac.pt, jleite@iscac.pt, maryam.abbasi@esg.ipsantarem.pt}
}

\keywords{Amyotrophic Lateral Sclerosis, Machine Learning, Class Imbalance,
Ensemble Learning, LightGBM, Bayesian Optimisation, SHAP, Explainability.}

\abstract{%
Amyotrophic lateral sclerosis (ALS) is a progressive and fatal neurodegenerative
disease in which a subset of patients, termed short survivors, die within
24 months of symptom onset. Identifying this subgroup at diagnosis is
clinically relevant but is complicated by severe class imbalance: short
survivors typically represent fewer than 15\% of ALS cohorts, and standard
classifiers tend to overlook them in favour of the majority class.
This paper replicates and extends the ensemble-imbalance machine learning
framework of Papaiz et al.~\cite{Papaiz2024} using the publicly available
PRO-ACT clinical database. Reproducing their cross-validation protocol with
BalancedBagging confirms that the ensemble approach significantly improves
balanced accuracy in six of seven classifiers (Bonferroni-corrected Wilcoxon
test, $p \leq 0.01$), with the MLP achieving the highest ensemble balanced
accuracy of 0.832. The framework is then extended by replacing grid search with
Bayesian hyperparameter optimisation (Optuna) and introducing LightGBM as a
seventh classifier across ten imbalance-handling strategies. On a held-out
partition drawn from the same PRO-ACT trial population, LightGBM with Random
Oversampling reaches an AUC of 0.923 and identifies 30 of 35 short-survival
patients (sensitivity 0.857). SVM with SMOTEENN achieves the highest sensitivity
(0.886) at a lower specificity (0.767). The MLP, despite leading
cross-validation AUC (0.915), identifies only one of those 35 patients at the
default decision threshold (sensitivity 0.029), illustrating that high AUC does
not guarantee adequate minority-class detection when a fixed threshold is applied.
SHAP analysis of both models identifies diagnostic delay and ALSFRS-R functional
decline slopes as the dominant predictors. These results are based on an internal
held-out set and external validation remains necessary.
}

\onecolumn \maketitle \normalsize \setcounter{footnote}{0} \vfill

\newpage

% ==========================================================================
\section{Introduction}
\label{sec:intro}
% ==========================================================================

Amyotrophic lateral sclerosis (ALS) is a progressive neurodegenerative disorder
affecting both upper and lower motor neurons, leading to generalised muscle
paralysis and, ultimately, respiratory failure. With a median survival of two to
five years from symptom onset, it is among the most severe neurological
conditions encountered in clinical practice~\cite{Feldman2022}. Within the ALS
population, a clinically important subset, referred to as short survivors, die
within 24 months of onset and follow a markedly faster disease trajectory. Early
identification of this subgroup has direct therapeutic relevance: it enables the
timely initiation of palliative and supportive care, supports eligible patients'
prioritisation in clinical trials, and facilitates advance care
planning~\cite{Miller2012}.

Machine learning (ML) approaches have increasingly been applied to ALS prognosis
prediction, drawing on clinical and functional data collected across large
multicentre registries~\cite{Grollemund2019,Papaiz2022}. The PRO-ACT (Pooled
Resource Open-Access ALS Clinical Trials) database~\cite{PROACT}, the largest
publicly available ALS data source, has served as the primary benchmark,
providing longitudinal records for thousands of trial participants. Studies have
demonstrated that features derived from the ALS Functional Rating Scale Revised
(ALSFRS-R)~\cite{cedarbaum1999alsfrs}, vital signs, and demographic information
can discriminate between survival trajectories~\cite{Ramamoorthy2022,Halbersberg2023}.

A persistent challenge in all such work is the pronounced class imbalance between
short and non-short survivors: short-survival patients typically represent only
10--15\% of ALS cohorts. Without specific corrections, standard classifiers
gravitate towards the majority class, yielding high specificity at the cost of
clinically unacceptable sensitivity~\cite{chawla2002smote}. Addressing this
imbalance has motivated approaches ranging from synthetic oversampling (SMOTE and
its variants) to ensemble-imbalance methods such as
BalancedBagging~\cite{lemaitre2017}, which integrates random undersampling within
each bootstrap estimator.

Papaiz et al.~\cite{Papaiz2024} proposed a systematic evaluation of six
classical ML classifiers combined with BalancedBagging, using 23 hand-crafted
clinical features derived from the PRO-ACT database. That study (hereafter the \emph{reference study}) showed that the ensemble wrapper
consistently improved balanced accuracy, with the MLP emerging as the best
performer under repeated cross-validation. The work left open several questions,
however: hyperparameters were tuned by grid search, gradient-boosted classifiers
were not considered, and performance was assessed entirely within cross-validation
without a separate held-out test set---a gap that, as our results show, has
non-trivial consequences.

This paper makes the following contributions:
\begin{enumerate}
  \item A faithful replication of the reference study on the publicly available
    PRO-ACT extract, confirmed by Bonferroni-corrected statistical tests across
    all seven classifiers;
  \item An extension replacing grid search with Bayesian hyperparameter
    optimisation (Optuna~\cite{optuna2019}) and adding LightGBM~\cite{ke2017lightgbm}
    as a seventh classifier, evaluated across ten imbalance-handling strategies;
  \item A held-out test evaluation showing that the MLP's leading cross-validation
    AUC does not translate to adequate minority-class sensitivity at the default
    threshold, whereas LightGBM with Random Oversampling achieves the best overall
    balance and SVM with SMOTEENN the highest raw sensitivity;
  \item SHAP-based interpretability analysis~\cite{lundberg2017shap,lundberg2020treeshap} for both the replication and the extension models,
    identifying diagnostic delay and ALSFRS-R decline slopes as the dominant
    prognostic variables.
\end{enumerate}

Section~\ref{sec:related} reviews related work. Section~\ref{sec:methodology}
describes the data and experimental design. Results appear in
Section~\ref{sec:results}, discussion in Section~\ref{sec:discussion}, and
Section~\ref{sec:conclusion} concludes.


% ==========================================================================
\section{Related Work}
\label{sec:related}
% ==========================================================================

\subsection{Machine Learning for ALS Prognosis}

ML approaches to ALS progression modelling have multiplied over the past
decade, with the PRO-ACT database serving as the shared benchmark for most of
them. Early crowdsourced efforts by Küffner et al.~\cite{kueffner2019} showed
that ensemble models could predict ALSFRS-R slope at accuracy levels competitive
with specialist clinical benchmarks. More recent work has explored deep
learning architectures for continuous trajectory modelling, including recurrent
and multi-task networks applied to longitudinal ALSFRS-R
data~\cite{Pancotti2022,Qin2025}. Ramamoorthy et al.~\cite{Ramamoorthy2022}
proposed a Gaussian-process framework for identifying progression patterns from
sparse longitudinal records.

At the level of binary survival classification, Halbersberg and
Lerner~\cite{Halbersberg2023} compared several ML methods for ALS stratification,
finding that ensemble-based approaches consistently outperformed individual
classifiers. Anani et al.~\cite{Anani2023} combined feature selection with
interpretable modelling to isolate the most prognostically informative variables.
Papaiz et al.~\cite{Papaiz2022}, in a systematic review, pinpointed class
imbalance and limited dataset size as the two recurring obstacles to reliable
ALS prognosis. Wang et al.~\cite{Wang2025} surveyed explainable AI across
neurodegenerative diseases, noting the widening gap between model complexity
and clinical interpretability.

\subsection{Imbalance-Aware Learning and Ensemble Methods}

Class imbalance in medical classification has attracted substantial methodological
attention. Resampling strategies include synthetic oversampling methods like
SMOTE~\cite{chawla2002smote}, BorderlineSMOTE, and ADASYN, as well as
undersampling methods such as Random Undersampling (RUS) and NearMiss. Hybrid
approaches (SMOTETomek and SMOTEENN) combine both paradigms to remove noisy
boundary samples alongside oversampling. The \texttt{imbalanced-learn}
library~\cite{lemaitre2017} provides scikit-learn-compatible
implementations of all these strategies, including BalancedBagging, which
applies random undersampling internally within each bootstrap estimator.
Which technique works best is highly problem-dependent; there is no universally
superior strategy, which is why systematic cross-validated comparison across
multiple classifiers remains the most defensible selection approach.

\subsection{Bayesian Hyperparameter Optimisation}

Grid search, as used in the reference study, evaluates hyperparameters on a fixed
predefined lattice, which scales poorly with the number of parameters and may miss
near-optimal configurations that lie between grid points. Bayesian optimisation
addresses this by building a probabilistic surrogate model of the objective
function and directing new evaluations towards promising regions. Optuna~\cite{optuna2019}, built around the Tree-structured Parzen Estimator
(TPE), has gained wide adoption in this space; its define-by-run API and
automatic trial pruning make it practical for large search spaces.

\subsection{Explainability in Clinical ML}

SHAP (SHapley Additive exPlanations)~\cite{lundberg2017shap} frames feature
attribution as a cooperative game: each feature is credited for its marginal
contribution to a prediction, averaged across all possible feature orderings.
For tree-based models, TreeSHAP~\cite{lundberg2020treeshap} computes these values
exactly in polynomial time. KernelSHAP extends the approach to architectures such
as MLPs via sampling-based approximation, at higher computational cost. Applied
to ALS prognosis~\cite{Anani2023}, SHAP has consistently ranked functional decline
metrics as the dominant predictors, a result that clinicians find immediately
recognisable.


% ==========================================================================
\section{Methodology}
\label{sec:methodology}
% ==========================================================================

\subsection{Dataset Overview}

The PRO-ACT database~\cite{PROACT} contains longitudinal clinical records from
ALS patients enrolled across multiple controlled trials, covering ALSFRS-R
assessments, vital signs, forced vital capacity (FVC), medication history, and
demographic data. Following the feature engineering protocol of
Papaiz et al.~\cite{Papaiz2024}, we extracted 23 categorical and ordinal
features representing each patient's clinical status at or near the time of
diagnosis.

The binary target variable distinguishes short survivors (death $\leq$ 24 months
from symptom onset) from non-short survivors (death $>$ 24 months, or alive
beyond 24 months). Patients alive at the time of their last observation but
without 24 months of confirmed follow-up were treated as censored and excluded,
consistent with the reference study. Patients missing site-of-onset records were
also excluded.

Table~\ref{tab:dataset} compares the resulting dataset with the one reported in
the reference study. Our cohort comprises 1,502 patients, 23.6\% fewer than the
original, primarily because the publicly available PRO-ACT extract examined
here contains a substantially higher proportion of missing FVC records (62.3\%
after survival and site filtering) and sparser height data required for BMI
derivation. As a consequence, the class imbalance ratio is slightly higher (7.6
versus 6.9).

\begin{table}[ht!]
  \centering
  \caption{Dataset comparison between the reference study~\protect\cite{Papaiz2024} and this work.}
  \label{tab:dataset}
  \resizebox{\linewidth}{!}{%
  \begin{tabular}{lcc}
    \toprule
    \textbf{Metric} & \textbf{Reference Study} & \textbf{This Work} \\
    \midrule
    Total patients   & 1,967          & 1,502          \\
    Short (\%)       & 13.0\% (256)   & 11.6\% (174)   \\
    Non-Short (\%)   & 87.0\% (1,711) & 88.4\% (1,328) \\
    Imbalance ratio  & 6.9            & 7.6            \\
    Features         & 23             & 23             \\
    \bottomrule
  \end{tabular}%
  }
\end{table}

\subsection{Data Preprocessing}

Preprocessing mirrored the reference study's pipeline. Binary features---sex,
site of onset, Riluzole use, and gastrostomy---were encoded as 0/1 indicators.
Age at onset was discretised into five ordinal groups ($\leq$39, 40--49, 50--59,
60--69, $\geq$70). Diagnostic delay was categorised as short ($\leq$8 months),
average (9--18 months), or long ($\geq$19 months), a proxy for pre-diagnostic
disease tempo. FVC at diagnosis was classified as abnormal ($<$80\% predicted) or
normal. BMI was derived from body weight and height measurements closest to the
diagnosis visit and encoded into four standard categories. ALSFRS-R functional
decline slopes for each of the ten subscores (Q1--Q10) were computed as the
rate of change per month from symptom onset using linear regression over
available assessments, and binned as slow ($<$0.05), average (0.05--0.14), or
rapid ($\geq$0.14 points per month). Disease regions involved
(bulbar, upper limb, lower limb, respiratory) were derived from ALSFRS-R
subscores at diagnosis.

\subsection{Data Partitioning}

The dataset was divided into a training set (80\%, $n = 1{,}201$; 139 short,
1,062 non-short) and a held-out test set (20\%, $n = 301$; 35 short, 266
non-short) using stratified random sampling to preserve the class distribution in
both subsets (random seed 42). The test set was sealed until all model selection
and hyperparameter optimisation had been completed.

\subsection{Experimental Setup}
\label{subsec:setup}

\paragraph{Classifiers.}
Seven classifiers were evaluated: K-Nearest Neighbours (KNN), Naïve Bayes (NB),
Decision Tree (DT), Random Forest (RF), Support Vector Machine (SVM),
Multi-Layer Perceptron (MLP), and LightGBM~\cite{ke2017lightgbm}. The first six
correspond directly to those used in the reference study; LightGBM constitutes this
work's primary algorithmic extension.

\paragraph{Imbalance-handling strategies.}
Ten class-imbalance correction strategies were evaluated: no correction (None),
Random Oversampling (ROS), SMOTE~\cite{chawla2002smote}, Borderline-SMOTE,
ADASYN, Random Undersampling (RUS), NearMiss, TomekLinks, SMOTETomek, and
SMOTEENN, yielding 70 classifier--technique combinations in total.

\paragraph{Hyperparameter optimisation.}
Replacing the original grid search, we employed Optuna~\cite{optuna2019} with
the TPE sampler for each of the 70 combinations, running 50 evaluation trials per
combination. ROC-AUC was used as the optimisation criterion, averaged over
cross-validation folds, to maintain direct comparability with the reference
study's evaluation metric. This choice was deliberate: by optimising for the same
criterion as the reference study, we can attribute any divergence between
cross-validation AUC and held-out minority-class performance to properties of the
criterion itself---rather than to differences in model selection procedure.
As Section~\ref{sec:results} demonstrates, this design decision brings the
AUC--sensitivity decoupling into sharp relief. Classifier-specific hyperparameter
search spaces were defined to cover a broad range of architectures (e.g.\ number
of layers and neurons for MLP; number of estimators and tree depth for RF and
LightGBM).

\paragraph{Cross-validation.}
Model selection was conducted using Repeated Stratified $K$-Fold
cross-validation with $K = 5$ and 3 repeats (15 evaluations per configuration),
applied exclusively to the training set.

\paragraph{BalancedBagging replication.}
To directly replicate the reference study, all seven classifiers were additionally
evaluated in two scenarios using the same cross-validation protocol: (i) a
single-model baseline with no imbalance correction, and (ii) wrapped inside a
BalancedBagging ensemble~\cite{lemaitre2017} with ten estimators and random
undersampling applied within each bootstrap iteration.

\paragraph{Statistical validation.}
Per-fold balanced accuracy scores from the 15 cross-validation folds were compared
between the single-model and BalancedBagging configurations for each classifier
using the Wilcoxon signed-rank test~\cite{wilcoxon1945}. Family-wise error rate
inflation across the seven simultaneous comparisons was controlled by applying
Bonferroni correction (corrected significance threshold
$\alpha_c = 0.05 / 7 \approx 0.007$).

\paragraph{SHAP interpretability.}
Feature importance was assessed using TreeSHAP~\cite{lundberg2020treeshap} for
LightGBM and KernelSHAP~\cite{lundberg2017shap} for the MLP (100 background
samples, 200 evaluation samples per instance). SHAP values were computed on the
held-out test set for both models after final retraining on the full training set.

\paragraph{Evaluation metrics.}
Models were assessed using ROC-AUC, balanced accuracy, sensitivity (recall for
the short-survival class), specificity, and geometric mean (G-Mean). Given that
failing to identify a short-survival patient constitutes a clinically serious
outcome, sensitivity received primary clinical emphasis throughout.


% ==========================================================================
\section{Results}
\label{sec:results}
% ==========================================================================

\subsection{Cross-Validation Performance}

Table~\ref{tab:cv} shows, for each classifier, the combination (out of 70
Optuna-optimised configurations) that achieved the highest AUC, averaged over
15 cross-validation folds.

\begin{table}[ht!]
  \centering
  \caption{Best cross-validation results per classifier (Optuna, 70 combinations,
    15-fold repeated CV). Sensitivity and specificity are for the short-survival
    (positive) class. G-Mean $= \sqrt{\text{Sensitivity} \times \text{Specificity}}$.}
  \label{tab:cv}
  \resizebox{\linewidth}{!}{%
  \begin{tabular}{llcccc}
    \toprule
    \textbf{Classifier} & \textbf{Best Technique} & \textbf{AUC}
      & \textbf{Sensitivity} & \textbf{Specificity} & \textbf{G-Mean} \\
    \midrule
    MLP           & None       & 0.915 & 0.427 & 0.969 & 0.641 \\
    LightGBM      & ROS        & 0.914 & 0.780 & 0.864 & 0.819 \\
    SVM           & SMOTEENN   & 0.911 & 0.904 & 0.784 & 0.842 \\
    Random Forest & None       & 0.906 & 0.276 & 0.983 & 0.517 \\
    Naïve Bayes   & None       & 0.891 & 0.727 & 0.867 & 0.793 \\
    KNN           & RUS        & 0.880 & 0.767 & 0.831 & 0.797 \\
    Decision Tree & TomekLinks & 0.861 & 0.413 & 0.948 & 0.621 \\
    \bottomrule
  \end{tabular}%
  }
\end{table}

The MLP achieves the highest cross-validation AUC (0.915) but at the expense of
low minority-class sensitivity (0.427). The Optuna configuration behind this
result is instructive: the best MLP carries a comparatively strong L2 penalty
($\alpha = 0.061$), roughly five times larger than the values found for the same
architecture under any resampling strategy. That regularisation pressure
suppresses the model's tendency to assign uniformly high scores to non-short
patients, improving probability ranking (and therefore AUC) without pushing
minority-class scores above the 0.5 decision boundary. The result is a model
that looks competitive in cross-validation but has already traded sensitivity for
ranking fidelity---a cross-validation echo of the same decoupling that becomes
fully visible at test time.
LightGBM with Random Oversampling delivers a comparable AUC (0.914) while
offering substantially better sensitivity (0.780) and G-Mean (0.819). Three
classifiers---MLP, Random Forest, and Naïve Bayes---attain their best AUC
without any imbalance correction, indicating that resampling does not improve
the model's overall ranking ability as measured by AUC for these architectures.
However, this masks an important limitation: for MLP and Random Forest in
particular, the high AUC coexists with low sensitivity, a pattern that the
results on the held-out test set will render unmistakable.

\subsection{Held-Out Test Set Evaluation}

After model selection was finalised on the training set, each classifier was
retrained on the full training partition using its best hyperparameter
configuration and evaluated on the sealed test set (35 short-survival,
266 non-short-survival patients). Results are presented in Table~\ref{tab:test}.

\begin{table}[ht!]
  \centering
  \caption{Held-out test set performance for each classifier under its best
    cross-validation configuration. Bold indicates the best value per metric.}
  \label{tab:test}
  \resizebox{\linewidth}{!}{%
  \begin{tabular}{llccccc}
    \toprule
    \textbf{Classifier} & \textbf{Technique} & \textbf{AUC}
      & \textbf{Bal.\ Acc.} & \textbf{Sensitivity} & \textbf{Specificity}
      & \textbf{G-Mean} \\
    \midrule
    LightGBM      & ROS        & \textbf{0.923} & \textbf{0.844} & 0.857
      & 0.831 & \textbf{0.844} \\
    Random Forest & None       & 0.919 & 0.623 & 0.257 & 0.989 & 0.504 \\
    MLP           & None       & 0.918 & 0.514 & 0.029 & \textbf{1.000} & 0.169 \\
    SVM           & SMOTEENN   & 0.909 & 0.826 & \textbf{0.886} & 0.767 & 0.824 \\
    Naïve Bayes   & None       & 0.889 & 0.774 & 0.714 & 0.835 & 0.772 \\
    KNN           & RUS        & 0.886 & 0.812 & 0.800 & 0.823 & 0.812 \\
    Decision Tree & TomekLinks & 0.874 & 0.628 & 0.286 & 0.970 & 0.526 \\
    \bottomrule
  \end{tabular}%
  }
\end{table}

The most consequential finding concerns the MLP. Despite ranking first in
cross-validation AUC, it detects only \textbf{one} of 35 short-survival patients
on the test set (sensitivity $= 0.029$, balanced accuracy $= 0.514$). This gap
between cross-validation ranking and held-out minority-class performance is
visualised in Figure~\ref{fig:cvtest}: it is a direct instance of the
\emph{AUC--sensitivity decoupling} effect under severe class imbalance, where a
model that correctly ranks probabilities on average may still compress
minority-class predicted probabilities below the decision threshold, rendering
them practically invisible at test time.

In contrast, LightGBM with Random Oversampling correctly identifies 30 of 35
short-survival patients (sensitivity $= 0.857$), while maintaining a specificity
of 0.831—effectively recovering 85.7\% of the patients who would die within
24 months. SVM with SMOTEENN achieves the highest raw sensitivity (0.886), though at a
lower specificity (0.767) and AUC (0.909). The two models serve different
priorities: SVM+SMOTEENN may be preferred where minimising false negatives is
the overriding concern, while LightGBM+ROS offers a stronger overall balance
across all metrics and the highest AUC (0.923).

\begin{figure}[ht!]
  \centering
  \includegraphics[width=0.95\linewidth]{figures/step6_cv_vs_test.png}
  \caption{Cross-validation AUC versus held-out test AUC for all seven
    classifiers. The MLP is the starkest case: it ranks first in CV AUC (0.915)
    yet achieves near-chance sensitivity on the test set (0.029), illustrating
    the AUC--sensitivity decoupling under severe class imbalance.}
  \label{fig:cvtest}
\end{figure}

\subsection{BalancedBagging Replication and Statistical Validation}

To replicate the reference study, all seven classifiers were evaluated in
single-model and BalancedBagging configurations under the same cross-validation
protocol. Table~\ref{tab:bonferroni} reports mean balanced accuracy for each
condition along with the statistical comparison.

\begin{table}[ht!]
  \centering
  \caption{Cross-validation balanced accuracy under single-model and
    BalancedBagging configurations, with per-fold Wilcoxon signed-rank tests
    (Bonferroni-corrected for 7 comparisons).
    $^{***}p \leq 0.001$; $^{**}p \leq 0.01$; n.s.\ = not significant.}
  \label{tab:bonferroni}
  \resizebox{\linewidth}{!}{%
  \begin{tabular}{lcccc}
    \toprule
    \textbf{Classifier} & \textbf{Single-Model} & \textbf{BalancedBagging}
      & \textbf{$\Delta$} & \textbf{$p$ (Bonf.)} \\
    \midrule
    KNN           & 0.566 & 0.788 & $+$0.221 & $<0.001^{***}$ \\
    Naïve Bayes   & 0.797 & 0.809 & $+$0.012 & 1.000 n.s.     \\
    Decision Tree & 0.653 & 0.777 & $+$0.125 & $0.001^{**}$   \\
    Random Forest & 0.629 & 0.814 & $+$0.185 & $<0.001^{***}$ \\
    SVM           & 0.500 & 0.793 & $+$0.293 & $<0.001^{***}$ \\
    MLP           & 0.698 & 0.832 & $+$0.134 & $<0.001^{***}$ \\
    LightGBM      & 0.692 & 0.820 & $+$0.128 & $<0.001^{***}$ \\
    \bottomrule
  \end{tabular}%
  }
\end{table}

Six of the seven classifiers exhibit a statistically significant improvement in
balanced accuracy under BalancedBagging ($p \leq 0.01$ after Bonferroni
correction). The sole exception is Naïve Bayes ($p = 1.0$, n.s.): its
generative modelling of class-conditional densities makes it relatively
insensitive to training class proportions, and its single-model balanced accuracy
(0.797) already approaches its ensemble counterpart (0.809). The SVM shows the
largest absolute gain ($\Delta = +0.293$), crossing from a balanced accuracy of
0.500 (equivalent to a classifier that predicts all instances as non-short) to
0.793 once internal balancing is applied within each bootstrap sample. The MLP
achieves the highest ensemble balanced accuracy across all classifiers (0.832),
consistent with the reference study's findings. LightGBM also improves
significantly under BalancedBagging ($\Delta = +0.128$, $p < 0.001$),
demonstrating that the ensemble strategy generalises beyond the six classifiers
considered in the original work.

\subsection{Comparison with the Reference Study}

Table~\ref{tab:comparison} summarises the methodological and performance differences
between the reference study and this work.

\begin{table}[ht!]
  \centering
  \caption{Direct comparison between Papaiz et al.~\cite{Papaiz2024} and this work.
    \textsuperscript{a}Best ensemble balanced accuracy reported under repeated CV.
    \textsuperscript{b}Not reported in the reference study.}
  \label{tab:comparison}
  \resizebox{\linewidth}{!}{%
  \begin{tabular}{lcc}
    \toprule
    \textbf{Aspect} & \textbf{Papaiz et al.~\cite{Papaiz2024}} & \textbf{This Work} \\
    \midrule
    Patients                    & 1,967              & 1,502              \\
    Classifiers evaluated       & 6                  & 7 (+LightGBM)      \\
    Imbalance strategies        & BalancedBagging    & 10 strategies + BB \\
    Hyperparameter optimisation & Grid search        & Bayesian (Optuna)  \\
    Best CV Bal.\ Acc.\ (MLP+BB)\textsuperscript{a} & $\approx$0.84 & 0.832 \\
    Held-out test evaluation    & No\textsuperscript{b} & Yes             \\
    Best test AUC               & —\textsuperscript{b} & 0.923 (LightGBM+ROS) \\
    Best test Sensitivity       & —\textsuperscript{b} & 0.857 (LightGBM+ROS) \\
    Explainability (SHAP)       & Yes (6 classifiers) & Yes (TreeSHAP + KernelSHAP) \\
    \bottomrule
  \end{tabular}%
  }
\end{table}

The CV balanced accuracy for MLP+BalancedBagging (0.832) comes close to the
reference figure ($\approx$0.84) despite a smaller and more imbalanced cohort,
which speaks to the consistency of the ensemble approach across different PRO-ACT
extracts. Beyond replication, three additions change the picture: Bayesian
optimisation surfaces hyperparameter configurations (notably the MLP's high L2
penalty) that grid search would rarely visit; LightGBM brings a gradient-boosted
option the original study did not include; and the held-out test reveals that the
MLP's cross-validation lead does not survive contact with unseen data.

\subsection{Feature Importance via SHAP}

Figure~\ref{fig:shap} shows the TreeSHAP output for LightGBM+ROS on the
held-out test set, with features ranked by mean absolute SHAP value.

\begin{figure}[ht!]
  \centering
  \includegraphics[width=0.95\linewidth]{figures/step7_lightgbm_fig3_combined.png}
  \caption{SHAP feature importance for LightGBM+ROS (held-out test set).
    \textit{Left:} global bar chart of mean $|\text{SHAP}|$ values.
    \textit{Right:} beeswarm plot showing per-instance SHAP contributions,
    coloured by feature value (red = high, blue = low).}
  \label{fig:shap}
\end{figure}

\texttt{Diagnosis\_Delay} is the single most influential predictor by a
substantial margin (mean $|\text{SHAP}| = 0.986$), more than double the value of
the next-ranked feature. A long diagnostic delay (typically indicating a
protracted, relatively benign pre-diagnostic period) is strongly associated with
non-short survival; conversely, a short delay signals rapid pre-diagnosis
progression and is a marker of short survival. A short diagnostic delay has long been recognised in the clinical literature as a marker of rapid pre-symptomatic progression~\cite{Feldman2022}.

ALSFRS-R functional decline slopes occupy positions two through four and six
in the global ranking, with Q8 (walking), Q4 (handwriting), Q1 (speech), and
Q10 (respiratory function) contributing substantially. Age at onset ranks fifth (mean $|\text{SHAP}| = 0.364$): older patients
generally progress faster, a relationship well documented in the ALS literature.
FVC at diagnosis appears at rank nine (mean $|\text{SHAP}| = 0.260$), adding a
respiratory dimension to the prognostic picture. The overall ranking closely
matches the SHAP results reported in the reference study~\cite{Papaiz2024}.

A parallel KernelSHAP analysis of the MLP identified the same top feature
(\texttt{Diagnosis\_Delay}) and a broadly similar ranking of ALSFRS-R slopes,
though with considerably smaller absolute attributions, a consequence of the
more diffuse weight distribution characteristic of neural network architectures.


% ==========================================================================
\section{Discussion}
\label{sec:discussion}
% ==========================================================================

\subsection{Replication Fidelity}

Our replication broadly confirms the central conclusion of the reference
study~\cite{Papaiz2024}: BalancedBagging consistently and significantly improves
classifier performance under repeated stratified cross-validation, with the MLP
attaining the best ensemble balanced accuracy (0.832 in this work versus
approximately 0.84 reported in the original). The slightly lower absolute values reflect the smaller cohort (1,502 versus
1,967 patients), the higher imbalance ratio (7.6 versus 6.9), and minor
differences in the PRO-ACT extract available to us. Across all classifiers,
the pattern of gains mirrors what the reference study reported, suggesting its
conclusions hold across different PRO-ACT extracts.

\subsection{The MLP Generalisation Problem}

The most consequential gap between cross-validation ranking and test performance
belongs to the MLP. Cross-validation ranked it first in AUC (0.915), yet it
identified only one of 35 short-survival patients at test time (sensitivity
0.029), despite an AUC of 0.918 on the same data. The high AUC confirms that
the model assigns higher scores to short-survival patients on average; the
near-zero sensitivity points to a threshold mismatch rather than an absence of
discriminative signal.

This pattern has a well-understood cause. AUC measures ranking ability---the
probability that a positive instance receives a higher predicted score than a
negative one---not the ability to classify correctly at a given threshold. Under
severe class imbalance, a model can learn to assign slightly higher probabilities
to short-survival patients on average while still placing those probabilities well
below the 0.5 decision boundary. Cross-validation detects the ranking signal but does not expose the threshold
mismatch. Lowering the decision threshold would likely recover sensitivity for
the MLP, at a cost to specificity; threshold calibration is a natural next step
that the current evaluation, conducted at the default 0.5 threshold, did not
include.

For practitioners, the implication is direct: in the presence of severe class
imbalance, AUC alone does not reliably reflect minority-class detection.
Metrics that directly penalise minority-class errors---sensitivity, balanced
accuracy, G-Mean---should be part of both the optimisation objective and the
final evaluation, and a held-out test is not optional.

\subsection{Advantages of LightGBM and Bayesian Optimisation}

LightGBM with Random Oversampling achieves the strongest test-set results across
all metrics: AUC 0.923, balanced accuracy 0.844, and sensitivity 0.857, correctly
identifying 30 of 35 short-survival patients. SVM with SMOTEENN reaches the
highest raw sensitivity (0.886) but at a lower specificity (0.767) and AUC
(0.909); it represents a viable alternative when minimising false negatives is
the sole clinical priority. Histogram-based gradient boosting corrects residuals
sequentially, allowing the model to progressively focus on hard-to-classify
minority instances rather than fitting a global decision boundary from a
disproportionate training distribution. Random Oversampling, by rebalancing the
training set before fitting, reinforces this tendency.

Bayesian optimisation (Optuna, TPE sampler, 50 trials) contributes a further edge. Unlike grid search, which covers a predefined lattice of values, TPE allocates
more trials to promising regions of the hyperparameter space as information
accumulates, making it particularly well-suited to the larger search spaces of
LightGBM compared to simpler classifiers.

Finally, TreeSHAP provides exact Shapley values for LightGBM in polynomial
time, a meaningful practical difference from the KernelSHAP approximation
required for the MLP, which demands considerably more computation and is
susceptible to sampling variance. In a clinical setting where individual
predictions need to be explained, this difference translates directly to the clinic.

\subsection{Limitations}

The available PRO-ACT extract yields 1,502 patients, 23.6\% fewer than the
original study, primarily due to missing FVC records, which constrains
statistical power and may inflate variance in cross-validation estimates.
With an imbalance ratio of 7.6, minority-class metrics are also sensitive to
small changes in the confusion matrix: a shift of two or three predictions can
move sensitivity by several percentage points. PRO-ACT participants were enrolled in controlled clinical trials and may differ
from community ALS patients in disease stage, comorbidity profile, and duration
of follow-up. Critically, the held-out test set is drawn from the same PRO-ACT
trial population as the training data; it provides an unbiased estimate of
in-distribution performance but does not replace external validation on independent
cohorts, which remains a necessary condition before these models could inform
clinical decisions. The feature set, while faithfully
reproduced from the reference study, represents each patient at a single time
point through categorical and ordinal encodings; richer longitudinal
representations and biomarker data are not captured.


% ==========================================================================
\section{Conclusion}
\label{sec:conclusion}
% ==========================================================================

This work replicates the ensemble-imbalance framework of Papaiz
et al.~\cite{Papaiz2024} and extends it with Bayesian hyperparameter
optimisation and LightGBM. The replication holds: across six of seven
classifiers, BalancedBagging produces statistically significant gains in
balanced accuracy ($p \leq 0.01$, Bonferroni-corrected), and the MLP reaches an
ensemble balanced accuracy of 0.832, consistent with the original study.

The held-out evaluation, drawn from the same PRO-ACT trial population, adds
three observations. First, the MLP's leading cross-validation AUC (0.915) does
not translate to minority-class detection at the default threshold: sensitivity
falls to 0.029, which reflects an AUC--sensitivity decoupling rather than an
absence of discriminative signal. Second, LightGBM with Random Oversampling
achieves the best overall balance across metrics (AUC 0.923, sensitivity 0.857,
balanced accuracy 0.844); SVM with SMOTEENN reaches higher raw sensitivity
(0.886) at lower specificity (0.767) and is a reasonable alternative when
minimising false negatives is the primary concern. Third, SHAP analysis of both
models consistently identifies diagnostic delay and ALSFRS-R decline slopes as
the dominant predictors.
Three methodological points follow from these results: AUC alone is insufficient
to assess minority-class detection under severe imbalance; threshold selection
should be part of any evaluation rather than left at 0.5 by default; and a
held-out test, even when internal, prevents overestimating performance based on
cross-validation rankings alone. External validation on cohorts independent of
PRO-ACT remains necessary before these findings can be applied in practice.
Future work should address threshold calibration and cost-sensitive learning,
richer longitudinal feature representations, and validation on independent
cohorts.


\bibliographystyle{apalike}
{\small
\bibliography{bibfile}}

\end{document}